\documentclass[11pt,a4paper]{article}
\usepackage[utf8]{inputenc}
\usepackage[T1]{fontenc}
\usepackage[spanish,es-minimal,es-noshorthands]{babel}
\usepackage{lmodern}
\usepackage[sfdefault,lining]{FiraSans}
\usepackage[a4paper,top=2.3cm,bottom=2.2cm,left=1.8cm,right=1.8cm,headheight=1.3cm,headsep=0.4cm,footskip=1cm]{geometry}
\usepackage{amsmath,amssymb}
\usepackage{icomma}
\usepackage{graphicx}
\usepackage[table]{xcolor}
\usepackage{tikz}
\usetikzlibrary{calc}
\usepackage{fontawesome5}
\usepackage{ccicons}
\usepackage{fancyhdr}
\usepackage{lastpage}
\usepackage{tasks}
\usepackage{needspace}
\usepackage[most]{tcolorbox}
\usepackage{array}
\usepackage[hidelinks]{hyperref}
\AtBeginDocument{\hypersetup{pdftitle={Ficha de ejercicios de la unidad 3: Programación lineal. 2º Bach.},pdfsubject={Matemáticas CCSS 2º Bach.},pdfkeywords={programación lineal, inecuaciones, sistemas de inecuaciones, región factible, vértices, función objetivo, rectas de nivel, PAU}}}

\definecolor{navy}{HTML}{1B2748}
\definecolor{accent}{HTML}{E8603A}
\definecolor{teal}{HTML}{1F9E90}
\definecolor{soft}{HTML}{EEF1F8}
\definecolor{muted}{HTML}{5D6685}

\setlength{\parindent}{0pt}
\setlength{\parskip}{3pt}
\settasks{label=\alph*), label-format=\bfseries\color{navy}, label-width=1.3em, item-indent=1.8em, column-sep=0.6em, after-item-skip=5pt, before-skip=3pt, after-skip=4pt}

% ---------- capçalera i peu ----------
\newcommand{\llicencia}{\href{https://creativecommons.org/licenses/by-nc-sa/4.0/deed.es}{\ccbyncsa}\ \ Licencia Creative Commons Atribución-NoComercial-CompartirIgual 4.0 Internacional (CC BY-NC-SA 4.0)}
\pagestyle{fancy}
\fancyhf{}
\fancyhead[L]{\raisebox{-0.35\height}{\includegraphics[height=0.95cm]{logo-gva.png}}}
\fancyhead[C]{\small\color{muted}Ficha de ejercicios · Unidad 3: Programación lineal · 2\textsuperscript{o} Bach.}
\fancyhead[R]{\raisebox{-0.35\height}{\includegraphics[height=0.95cm]{logo-vermellar.png}}}
\fancyfoot[L]{\scriptsize\color{muted}\llicencia}
\fancyfoot[R]{\small\color{muted}\thepage}
\renewcommand{\headrulewidth}{0.4pt}
\renewcommand{\headrule}{\color{muted!40}\hrule height \headrulewidth}
\fancypagestyle{primera}{\fancyhf{}\renewcommand{\headrulewidth}{0pt}
  \fancyfoot[L]{\scriptsize\color{muted}\llicencia}
  \fancyfoot[R]{\small\color{muted}\thepage}}

% ---------- nivells i exercicis ----------
\newcommand{\estrelles}[1]{{\color{accent}\ifcase#1\or\faStar\faStar[regular]\faStar[regular]\or\faStar\faStar\faStar[regular]\or\faStar\faStar\faStar\fi}}
\newcounter{exer}
\newcommand{\ex}[1]{\par\needspace{4\baselineskip}\medskip\refstepcounter{exer}\noindent{\color{navy}\bfseries\theexer.}\ \ }

\newcommand{\seccio}[2]{\par\needspace{10\baselineskip}\bigskip
  \begin{tcolorbox}[colback=navy,colframe=navy,coltext=white,arc=2mm,boxrule=0pt,left=3mm,right=3mm,top=1.2mm,bottom=1.2mm]
  {\large\bfseries #1\quad #2}
  \end{tcolorbox}\vspace{-2pt}}


% quadre de recordatori (només a la fitxa d'exercicis)
\newcommand{\recorda}[1]{\par\smallskip\begin{tcolorbox}[colback=accent!7,colframe=accent!55,arc=2mm,boxrule=0.5pt,left=2.5mm,right=2.5mm,top=1mm,bottom=1mm]\small{\color{accent}\faLightbulb[regular]}\ \textbf{Recuerda:}\ #1\end{tcolorbox}}
% etiqueta dels exercicis de la PAU
\newcommand{\pau}[1]{{\color{teal}\small\bfseries\faUniversity\ PAU · #1.}\ }
\newcommand{\res}[1]{#1}\newcommand{\rest}[1]{#1}

% =====================================================================
% Capçalera FentMates de les fitxes: un color per a cada curs.
% Els colors són els de la pàgina inicial (LEVELS a assets/js/materials.js);
% el logo ve de _eines/logo/pdf/ (python3 _eines/logo/genera.py).
%
% Ús, al final del preàmbul de la unitat:
%   % =====================================================================
% Capçalera FentMates de les fitxes: un color per a cada curs.
% Els colors són els de la pàgina inicial (LEVELS a assets/js/materials.js);
% el logo ve de _eines/logo/pdf/ (python3 _eines/logo/genera.py).
%
% Ús, al final del preàmbul de la unitat:
%   % =====================================================================
% Capçalera FentMates de les fitxes: un color per a cada curs.
% Els colors són els de la pàgina inicial (LEVELS a assets/js/materials.js);
% el logo ve de _eines/logo/pdf/ (python3 _eines/logo/genera.py).
%
% Ús, al final del preàmbul de la unitat:
%   \input{../comu/fentmates.tex}
%   \fmcurs{1eso}
% i, al començament del document:
%   \capcalera{Matemàtiques}{1r ESO}{Fitxa d'exercicis · Unitat 2}{Divisibilitat}
% Opcional abans de \input: \def\fmneed{14} (espai mínim abans de cada secció).
% =====================================================================
\makeatletter
\def\fm@logodir{../../logo/pdf}
\providecommand{\fmneed}{10}
\def\fm@cursnom{}

% clau del curs, color i color del text damunt del color
\newcommand{\fm@curs}[3]{\@namedef{fm@color@#1}{#2}\@namedef{fm@text@#1}{#3}}
\fm@curs{1eso}{E6B84A}{navy}
\fm@curs{2eso}{3B82F6}{white}
\fm@curs{3eso}{FF7A50}{white}
\fm@curs{4eso-a}{A78BFA}{white}
\fm@curs{4eso-b}{6D28D9}{white}
\fm@curs{1batx-mi}{9B1C1C}{white}
\fm@curs{1batx-ccss}{EF6B6B}{white}
\fm@curs{2batx-mii}{1E6B47}{white}
\fm@curs{2batx-ccss}{2BB3A3}{white}

\newcommand{\fmcurs}[1]{%
  \@ifundefined{fm@color@#1}{\PackageError{fentmates}{Curs desconegut: #1}{}}{}%
  \edef\fm@tmp{\noexpand\definecolor{curs}{HTML}{\@nameuse{fm@color@#1}}%
               \noexpand\colorlet{curstext}{\@nameuse{fm@text@#1}}}\fm@tmp
  \colorlet{cursfosc}{curs!72!black}%
  \def\fm@logo{\fm@logodir/vertical-#1.pdf}%
  \def\fm@marca{\fm@logodir/marca-#1.pdf}}

% ---------- capçalera de la primera pàgina ----------
% #1 matèria (petit), #2 curs, #3 tipus de document i unitat, #4 títol
\newcommand{\capcalera}[4]{%
  \gdef\fm@cursnom{#2}%
  \par\noindent
  \begin{tcolorbox}[colback=curs,colframe=curs,boxrule=0pt,arc=2.2mm,width=3.5cm,height=2.15cm,
      valign=center,left=3mm,right=2mm,top=1mm,bottom=1mm,nobeforeafter,box align=center]
    \color{curstext}{\footnotesize #1}\par\vspace{1pt}{\fontsize{19}{21}\selectfont\bfseries #2\par}%
  \end{tcolorbox}\hspace{5mm}%
  \begin{minipage}[c]{\dimexpr\linewidth-3.5cm-5mm-2.9cm\relax}
    {\color{cursfosc}\small\bfseries\MakeUppercase{#3}}\par\vspace{2pt}
    {\color{navy}\fontsize{20}{23}\selectfont\bfseries #4\par}
  \end{minipage}\hfill
  \begin{minipage}[c]{2.6cm}\raggedleft\includegraphics[width=2.6cm]{\fm@logo}\end{minipage}\par
  \vspace{2pt}}

% ---------- barres de secció del color del curs ----------
\renewcommand{\seccio}[2]{\par\needspace{\fmneed\baselineskip}\bigskip
  \begin{tcolorbox}[colback=curs,colframe=curs,coltext=curstext,arc=2mm,boxrule=0pt,left=3mm,right=3mm,top=1.2mm,bottom=1.2mm]
  {\large\bfseries #1\quad #2}
  \end{tcolorbox}\vspace{-2pt}}

% ---------- estrelles de dificultat del color del curs (a Batxillerat no se'n posen) ----------
\@ifundefined{estrelles}{}{%
  \renewcommand{\estrelles}[1]{{\color{curs}\ifcase#1\or\faStar\faStar[regular]\faStar[regular]\or\faStar\faStar\faStar[regular]\or\faStar\faStar\faStar\fi}}}

% ---------- capçalera i peu de les pàgines ----------
% primera pàgina: logos institucionals al peu
\fancypagestyle{primera}{\fancyhf{}\renewcommand{\headrulewidth}{0pt}
  \fancyfoot[L]{\raisebox{-0.32\height}{\includegraphics[height=0.78cm]{logo-gva.png}}\hspace{5mm}%
    \raisebox{-0.32\height}{\includegraphics[height=0.78cm]{logo-vermellar.png}}}
  \fancyfoot[R]{\scriptsize\color{muted}\href{https://creativecommons.org/licenses/by-nc-sa/4.0/deed.ca}{\ccbyncsa}\ \ CC BY-NC-SA 4.0 · fentmates.org\qquad\small\thepage}}
% la resta: el símbol a l'esquerra i el curs a la dreta, amb el títol al centre com abans
\fancyhead[L]{\raisebox{-0.3\height}{\includegraphics[height=0.62cm]{\fm@marca}}}
\fancyhead[R]{\small\bfseries\color{cursfosc}\fm@cursnom}
\renewcommand{\headrule}{\color{curs!45}\hrule height \headrulewidth}
\makeatother

%   \fmcurs{1eso}
% i, al començament del document:
%   \capcalera{Matemàtiques}{1r ESO}{Fitxa d'exercicis · Unitat 2}{Divisibilitat}
% Opcional abans de \input: \def\fmneed{14} (espai mínim abans de cada secció).
% =====================================================================
\makeatletter
\def\fm@logodir{../../logo/pdf}
\providecommand{\fmneed}{10}
\def\fm@cursnom{}

% clau del curs, color i color del text damunt del color
\newcommand{\fm@curs}[3]{\@namedef{fm@color@#1}{#2}\@namedef{fm@text@#1}{#3}}
\fm@curs{1eso}{E6B84A}{navy}
\fm@curs{2eso}{3B82F6}{white}
\fm@curs{3eso}{FF7A50}{white}
\fm@curs{4eso-a}{A78BFA}{white}
\fm@curs{4eso-b}{6D28D9}{white}
\fm@curs{1batx-mi}{9B1C1C}{white}
\fm@curs{1batx-ccss}{EF6B6B}{white}
\fm@curs{2batx-mii}{1E6B47}{white}
\fm@curs{2batx-ccss}{2BB3A3}{white}

\newcommand{\fmcurs}[1]{%
  \@ifundefined{fm@color@#1}{\PackageError{fentmates}{Curs desconegut: #1}{}}{}%
  \edef\fm@tmp{\noexpand\definecolor{curs}{HTML}{\@nameuse{fm@color@#1}}%
               \noexpand\colorlet{curstext}{\@nameuse{fm@text@#1}}}\fm@tmp
  \colorlet{cursfosc}{curs!72!black}%
  \def\fm@logo{\fm@logodir/vertical-#1.pdf}%
  \def\fm@marca{\fm@logodir/marca-#1.pdf}}

% ---------- capçalera de la primera pàgina ----------
% #1 matèria (petit), #2 curs, #3 tipus de document i unitat, #4 títol
\newcommand{\capcalera}[4]{%
  \gdef\fm@cursnom{#2}%
  \par\noindent
  \begin{tcolorbox}[colback=curs,colframe=curs,boxrule=0pt,arc=2.2mm,width=3.5cm,height=2.15cm,
      valign=center,left=3mm,right=2mm,top=1mm,bottom=1mm,nobeforeafter,box align=center]
    \color{curstext}{\footnotesize #1}\par\vspace{1pt}{\fontsize{19}{21}\selectfont\bfseries #2\par}%
  \end{tcolorbox}\hspace{5mm}%
  \begin{minipage}[c]{\dimexpr\linewidth-3.5cm-5mm-2.9cm\relax}
    {\color{cursfosc}\small\bfseries\MakeUppercase{#3}}\par\vspace{2pt}
    {\color{navy}\fontsize{20}{23}\selectfont\bfseries #4\par}
  \end{minipage}\hfill
  \begin{minipage}[c]{2.6cm}\raggedleft\includegraphics[width=2.6cm]{\fm@logo}\end{minipage}\par
  \vspace{2pt}}

% ---------- barres de secció del color del curs ----------
\renewcommand{\seccio}[2]{\par\needspace{\fmneed\baselineskip}\bigskip
  \begin{tcolorbox}[colback=curs,colframe=curs,coltext=curstext,arc=2mm,boxrule=0pt,left=3mm,right=3mm,top=1.2mm,bottom=1.2mm]
  {\large\bfseries #1\quad #2}
  \end{tcolorbox}\vspace{-2pt}}

% ---------- estrelles de dificultat del color del curs (a Batxillerat no se'n posen) ----------
\@ifundefined{estrelles}{}{%
  \renewcommand{\estrelles}[1]{{\color{curs}\ifcase#1\or\faStar\faStar[regular]\faStar[regular]\or\faStar\faStar\faStar[regular]\or\faStar\faStar\faStar\fi}}}

% ---------- capçalera i peu de les pàgines ----------
% primera pàgina: logos institucionals al peu
\fancypagestyle{primera}{\fancyhf{}\renewcommand{\headrulewidth}{0pt}
  \fancyfoot[L]{\raisebox{-0.32\height}{\includegraphics[height=0.78cm]{logo-gva.png}}\hspace{5mm}%
    \raisebox{-0.32\height}{\includegraphics[height=0.78cm]{logo-vermellar.png}}}
  \fancyfoot[R]{\scriptsize\color{muted}\href{https://creativecommons.org/licenses/by-nc-sa/4.0/deed.ca}{\ccbyncsa}\ \ CC BY-NC-SA 4.0 · fentmates.org\qquad\small\thepage}}
% la resta: el símbol a l'esquerra i el curs a la dreta, amb el títol al centre com abans
\fancyhead[L]{\raisebox{-0.3\height}{\includegraphics[height=0.62cm]{\fm@marca}}}
\fancyhead[R]{\small\bfseries\color{cursfosc}\fm@cursnom}
\renewcommand{\headrule}{\color{curs!45}\hrule height \headrulewidth}
\makeatother

%   \fmcurs{1eso}
% i, al començament del document:
%   \capcalera{Matemàtiques}{1r ESO}{Fitxa d'exercicis · Unitat 2}{Divisibilitat}
% Opcional abans de \input: \def\fmneed{14} (espai mínim abans de cada secció).
% =====================================================================
\makeatletter
\def\fm@logodir{../../logo/pdf}
\providecommand{\fmneed}{10}
\def\fm@cursnom{}

% clau del curs, color i color del text damunt del color
\newcommand{\fm@curs}[3]{\@namedef{fm@color@#1}{#2}\@namedef{fm@text@#1}{#3}}
\fm@curs{1eso}{E6B84A}{navy}
\fm@curs{2eso}{3B82F6}{white}
\fm@curs{3eso}{FF7A50}{white}
\fm@curs{4eso-a}{A78BFA}{white}
\fm@curs{4eso-b}{6D28D9}{white}
\fm@curs{1batx-mi}{9B1C1C}{white}
\fm@curs{1batx-ccss}{EF6B6B}{white}
\fm@curs{2batx-mii}{1E6B47}{white}
\fm@curs{2batx-ccss}{2BB3A3}{white}

\newcommand{\fmcurs}[1]{%
  \@ifundefined{fm@color@#1}{\PackageError{fentmates}{Curs desconegut: #1}{}}{}%
  \edef\fm@tmp{\noexpand\definecolor{curs}{HTML}{\@nameuse{fm@color@#1}}%
               \noexpand\colorlet{curstext}{\@nameuse{fm@text@#1}}}\fm@tmp
  \colorlet{cursfosc}{curs!72!black}%
  \def\fm@logo{\fm@logodir/vertical-#1.pdf}%
  \def\fm@marca{\fm@logodir/marca-#1.pdf}}

% ---------- capçalera de la primera pàgina ----------
% #1 matèria (petit), #2 curs, #3 tipus de document i unitat, #4 títol
\newcommand{\capcalera}[4]{%
  \gdef\fm@cursnom{#2}%
  \par\noindent
  \begin{tcolorbox}[colback=curs,colframe=curs,boxrule=0pt,arc=2.2mm,width=3.5cm,height=2.15cm,
      valign=center,left=3mm,right=2mm,top=1mm,bottom=1mm,nobeforeafter,box align=center]
    \color{curstext}{\footnotesize #1}\par\vspace{1pt}{\fontsize{19}{21}\selectfont\bfseries #2\par}%
  \end{tcolorbox}\hspace{5mm}%
  \begin{minipage}[c]{\dimexpr\linewidth-3.5cm-5mm-2.9cm\relax}
    {\color{cursfosc}\small\bfseries\MakeUppercase{#3}}\par\vspace{2pt}
    {\color{navy}\fontsize{20}{23}\selectfont\bfseries #4\par}
  \end{minipage}\hfill
  \begin{minipage}[c]{2.6cm}\raggedleft\includegraphics[width=2.6cm]{\fm@logo}\end{minipage}\par
  \vspace{2pt}}

% ---------- barres de secció del color del curs ----------
\renewcommand{\seccio}[2]{\par\needspace{\fmneed\baselineskip}\bigskip
  \begin{tcolorbox}[colback=curs,colframe=curs,coltext=curstext,arc=2mm,boxrule=0pt,left=3mm,right=3mm,top=1.2mm,bottom=1.2mm]
  {\large\bfseries #1\quad #2}
  \end{tcolorbox}\vspace{-2pt}}

% ---------- estrelles de dificultat del color del curs (a Batxillerat no se'n posen) ----------
\@ifundefined{estrelles}{}{%
  \renewcommand{\estrelles}[1]{{\color{curs}\ifcase#1\or\faStar\faStar[regular]\faStar[regular]\or\faStar\faStar\faStar[regular]\or\faStar\faStar\faStar\fi}}}

% ---------- capçalera i peu de les pàgines ----------
% primera pàgina: logos institucionals al peu
\fancypagestyle{primera}{\fancyhf{}\renewcommand{\headrulewidth}{0pt}
  \fancyfoot[L]{\raisebox{-0.32\height}{\includegraphics[height=0.78cm]{logo-gva.png}}\hspace{5mm}%
    \raisebox{-0.32\height}{\includegraphics[height=0.78cm]{logo-vermellar.png}}}
  \fancyfoot[R]{\scriptsize\color{muted}\href{https://creativecommons.org/licenses/by-nc-sa/4.0/deed.ca}{\ccbyncsa}\ \ CC BY-NC-SA 4.0 · fentmates.org\qquad\small\thepage}}
% la resta: el símbol a l'esquerra i el curs a la dreta, amb el títol al centre com abans
\fancyhead[L]{\raisebox{-0.3\height}{\includegraphics[height=0.62cm]{\fm@marca}}}
\fancyhead[R]{\small\bfseries\color{cursfosc}\fm@cursnom}
\renewcommand{\headrule}{\color{curs!45}\hrule height \headrulewidth}
\makeatother

\fmcurs{2batx-ccss}

\begin{document}
\thispagestyle{primera}

\capcalera{Matemáticas CCSS}{2\textsuperscript{o} Bach.}{Ficha de ejercicios · Unidad 3}{Programación lineal}
\vspace{10pt}
Nombre y apellidos: \rule{6.6cm}{0.4pt}\hfill Grupo: \rule{1.6cm}{0.4pt}\hfill Fecha: \rule{2.2cm}{0.4pt}
\vspace{8pt}

\begin{tcolorbox}[colback=soft,colframe=soft,arc=2mm,left=3mm,right=3mm,top=1.5mm,bottom=1.5mm]
\small Haz las gráficas con regla y papel cuadriculado, e indica siempre los vértices de la región. Los ejercicios marcados con {\color{teal}\bfseries\faUniversity\ PAU} son de las pruebas de acceso a la universidad de la Comunidad Valenciana; encontrarás muchos más, resueltos, en la página de ejercicios PAU de la unidad.
\end{tcolorbox}

% =====================================================
\seccio{1}{Inecuaciones lineales con dos incógnitas}

\ex{1} Averigua cuáles de los puntos siguientes son solución de la inecuación $3x-2y\leq4$:\begin{tasks}(4)
  \task $(0,\,0)$
  \task $(2,\,1)$
  \task $(-1,\,3)$
  \task $(4,\,-2)$
\end{tasks}\recorda{Un punto $(x_0,y_0)$ es solución si, al sustituirlo, la desigualdad es cierta.}

\ex{1} Resuelve gráficamente las inecuaciones siguientes:\begin{tasks}(4)
  \task $x+y\leq4$
  \task $2x-y>2$
  \task $y\geq3x$
  \task $x-3y<6$
\end{tasks}\recorda{1) Dibuja la recta frontera: continua si es $\leq,\geq$ y discontinua si es $<,>$. \ 2) Elige un punto de prueba que no sea de la recta (normalmente $(0,\,0)$). \ 3) Si lo cumple, la solución es el semiplano que lo contiene; si no, el otro.}

\ex{1} Resuelve gráficamente las inecuaciones siguientes, en las que solo aparece una de las incógnitas:\begin{tasks}(4)
  \task $x\geq-2$
  \task $y<3$
  \task $2x\leq5$
  \task $-y\geq1$
\end{tasks}\recorda{La recta $x=k$ es \textbf{vertical} e $y=k$ es \textbf{horizontal}. \ $x\geq k$: a la derecha; \ $x\leq k$: a la izquierda; \ $y\geq k$: por encima; \ $y\leq k$: por debajo. No hace falta punto de prueba.}

\ex{2} Escribe la inecuación que tiene como solución cada semiplano sombreado:\begin{tasks}(2)
  \task \begin{tikzpicture}[x=0.5250cm,y=0.5250cm,baseline=(current bounding box.center)]\begin{scope}\clip (-2,-2) rectangle (6,4);\foreach \i in {-2,0,...,6} \draw[muted!25,very thin] (\i,-2)--(\i,4);\foreach \i in {-2,-1,...,4} \draw[muted!25,very thin] (-2,\i)--(6,\i);\fill[teal!22] (-2.0000,-2.0000)--(6.0000,-2.0000)--(6.0000,-1.0000)--(-2.0000,3.0000)--cycle;\end{scope}\draw[-latex,thick] (-2,0)--(6.5,0) node[below left=-1pt,font=\scriptsize]{$x$};\draw[-latex,thick] (0,-2)--(0,4.25) node[below left=-1pt,font=\scriptsize]{$y$};\node[below,font=\tiny,text=muted] at (-2,0) {$-2$};\node[below,font=\tiny,text=muted] at (2,0) {$2$};\node[below,font=\tiny,text=muted] at (4,0) {$4$};\node[below,font=\tiny,text=muted] at (6,0) {$6$};\node[left,font=\tiny,text=muted] at (0,-2) {$-2$};\node[left,font=\tiny,text=muted] at (0,-1) {$-1$};\node[left,font=\tiny,text=muted] at (0,1) {$1$};\node[left,font=\tiny,text=muted] at (0,2) {$2$};\node[left,font=\tiny,text=muted] at (0,3) {$3$};\node[left,font=\tiny,text=muted] at (0,4) {$4$};\begin{scope}\clip (-2,-2) rectangle (6,4);\draw[accent,line width=1.1pt] (-3.0000,3.5000)--(7.0000,-1.5000);\end{scope}\end{tikzpicture}
  \task \begin{tikzpicture}[x=0.7000cm,y=0.5750cm,baseline=(current bounding box.center)]\begin{scope}\clip (-2,-3) rectangle (4,5);\foreach \i in {-2,-1,...,4} \draw[muted!25,very thin] (\i,-3)--(\i,5);\foreach \i in {-2,0,...,5} \draw[muted!25,very thin] (-2,\i)--(4,\i);\fill[teal!22] (-2.0000,-3.0000)--(-1.0000,-3.0000)--(3.0000,5.0000)--(-2.0000,5.0000)--cycle;\end{scope}\draw[-latex,thick] (-2,0)--(4.25,0) node[below left=-1pt,font=\scriptsize]{$x$};\draw[-latex,thick] (0,-3)--(0,5.5) node[below left=-1pt,font=\scriptsize]{$y$};\node[below,font=\tiny,text=muted] at (-2,0) {$-2$};\node[below,font=\tiny,text=muted] at (-1,0) {$-1$};\node[below,font=\tiny,text=muted] at (1,0) {$1$};\node[below,font=\tiny,text=muted] at (2,0) {$2$};\node[below,font=\tiny,text=muted] at (3,0) {$3$};\node[below,font=\tiny,text=muted] at (4,0) {$4$};\node[left,font=\tiny,text=muted] at (0,-2) {$-2$};\node[left,font=\tiny,text=muted] at (0,2) {$2$};\node[left,font=\tiny,text=muted] at (0,4) {$4$};\begin{scope}\clip (-2,-3) rectangle (4,5);\draw[accent,line width=1.1pt,dashed] (-1.5000,-4.0000)--(3.5000,6.0000);\end{scope}\end{tikzpicture}
\end{tasks}\recorda{Halla la ecuación de la recta a partir de dos puntos y decide el signo con un punto de la zona sombreada.}

\ex{2} Expresa cada frase con una inecuación y represéntala gráficamente ($x$ e $y$ son dos números reales):\begin{tasks}(1)
  \task El doble de un número más el triple de otro no supera 12.
  \task La diferencia entre el primer número y el segundo es, como mínimo, 2.
\end{tasks}


% =====================================================
\seccio{2}{Sistemas de inecuaciones. Región factible y vértices}

\ex{1} Resuelve gráficamente los sistemas siguientes y calcula los vértices de la región solución. Indica si es acotada o no:\begin{tasks}(2)
  \task $\left\{\begin{aligned}x+y&\leq5\\x-y&\geq-1\\x\geq0,\ y&\geq0\end{aligned}\right.$
  \task $\left\{\begin{aligned}2x+y&\geq4\\x+2y&\geq5\\x\geq0,\ y&\geq0\end{aligned}\right.$
\end{tasks}\recorda{La solución es la zona común a todos los semiplanos. Los vértices son los puntos de corte de dos rectas frontera que cumplen \textbf{todas} las inecuaciones.}

\ex{1} Representa la región definida por las inecuaciones siguientes y calcula sus vértices:\[1\leq x\leq4,\qquad 0\leq y\leq3,\qquad x+y\leq6\]\recorda{Una doble desigualdad $a\leq x\leq b$ son dos inecuaciones: $x\geq a$ y $x\leq b$ (una franja vertical).}

\ex{2} Resuelve gráficamente los sistemas siguientes:\begin{tasks}(2)
  \task $\left\{\begin{aligned}x+y&\geq2\\x+y&\leq1\\x\geq0,\ y&\geq0\end{aligned}\right.$
  \task $\left\{\begin{aligned}y&\geq x\\y&\leq x+2\\0\leq x&\leq3\end{aligned}\right.$
\end{tasks}

\ex{2} Escribe el sistema de inecuaciones que tiene como solución la región sombreada:

\begin{center}\begin{tikzpicture}[x=0.6250cm,y=0.6250cm,baseline=(current bounding box.center)]\begin{scope}\clip (-1,-1) rectangle (7,5.5);\foreach \i in {0,2,...,7} \draw[muted!25,very thin] (\i,-1)--(\i,5.5);\foreach \i in {-1,0,...,5.5} \draw[muted!25,very thin] (-1,\i)--(7,\i);\fill[teal!22] (6.0000,0.0000)--(2.0000,4.0000)--(0.0000,0.0000)--cycle;\end{scope}\draw[-latex,thick] (-1,0)--(7.5,0) node[below left=-1pt,font=\scriptsize]{$x$};\draw[-latex,thick] (0,-1)--(0,5.75) node[below left=-1pt,font=\scriptsize]{$y$};\node[below,font=\tiny,text=muted] at (2,0) {$2$};\node[below,font=\tiny,text=muted] at (4,0) {$4$};\node[below,font=\tiny,text=muted] at (6,0) {$6$};\node[left,font=\tiny,text=muted] at (0,-1) {$-1$};\node[left,font=\tiny,text=muted] at (0,1) {$1$};\node[left,font=\tiny,text=muted] at (0,2) {$2$};\node[left,font=\tiny,text=muted] at (0,3) {$3$};\node[left,font=\tiny,text=muted] at (0,4) {$4$};\node[left,font=\tiny,text=muted] at (0,5) {$5$};\begin{scope}\clip (-1,-1) rectangle (7,5.5);\draw[accent,line width=1.1pt] (-1.0000,-2.0000)--(3.2500,6.5000);\draw[blue!70!black,line width=1.1pt] (8.0000,-2.0000)--(-0.5000,6.5000);\end{scope}\filldraw[fill=white,draw=navy] (6.0000,0.0000) circle (1.5pt);\node[font=\tiny,inner sep=1.5pt,fill=white,fill opacity=.75,text opacity=1] at ($(6.0000,0.0000)+(-21.8:0.36cm)$) {$(6,\,0)$};\filldraw[fill=white,draw=navy] (2.0000,4.0000) circle (1.5pt);\node[font=\tiny,inner sep=1.5pt,fill=white,fill opacity=.75,text opacity=1] at ($(2.0000,4.0000)+(104.0:0.36cm)$) {$(2,\,4)$};\filldraw[fill=white,draw=navy] (0.0000,0.0000) circle (1.5pt);\node[font=\tiny,inner sep=1.5pt,fill=white,fill opacity=.75,text opacity=1] at ($(0.0000,0.0000)+(-153.4:0.36cm)$) {$(0,\,0)$};\end{tikzpicture}\end{center}

\ex{2} Considera el sistema $\ x+2y\leq8,\ \ 3x+y\leq9,\ \ x\geq0,\ \ y\geq0$.\begin{tasks}(1)
  \task Averigua si los puntos $(1,\,2)$, $(2,\,3)$ y $(3,\,1)$ pertenecen a la región factible.
  \task Representa la región y calcula sus vértices.
\end{tasks}

\ex{2} Representa la región definida por $\ x\geq1,\ \ y\geq2,\ \ x+y\geq5\ $. ¿Es acotada? Calcula sus vértices.


% =====================================================
\seccio{3}{Optimización de una función objetivo sobre una región}

\ex{1} Calcula el máximo y el mínimo de la función $f(x,y)=3x+2y$ en la región del apartado a) del ejercicio 6.\recorda{Si la región es acotada, el máximo y el mínimo se alcanzan en vértices: evalúa la función en todos.}

\ex{1} La región factible de un problema es el polígono de vértices $A(0,\,0)$, $B(6,\,0)$, $C(5,\,3)$, $D(2,\,5)$ y $E(0,\,4)$. Halla el máximo y el mínimo de las funciones $f(x,y)=x+2y$ y $g(x,y)=2x-y$ en esta región.

\ex{2} Calcula, si existen, el máximo y el mínimo de $f(x,y)=4x+3y$ en la región del apartado b) del ejercicio 6.\recorda{Si la región no es acotada, puede que no haya máximo (o mínimo). Compruébalo con las rectas de nivel $4x+3y=k$.}

\ex{2} Resuelve el problema anterior gráficamente con rectas de nivel: dibuja la recta $4x+3y=0$ y desplázala paralelamente. ¿Cuál es el primer vértice de la región que toca?


% =====================================================
\seccio{4}{Problemas de programación lineal}

\ex{1} Una carpintería fabrica mesas y sillas. Cada mesa necesita 3 horas de montaje y 1 hora de acabado, y cada silla, 1 hora de montaje y 1 hora de acabado. Cada semana se dispone de 60 horas de montaje y 30 horas de acabado. El beneficio es de 40 € por mesa y de 20 € por silla. ¿Cuántas mesas y cuántas sillas hay que fabricar para obtener el máximo beneficio?\recorda{Pasos: variables (tabla de datos) $\to$ función objetivo $\to$ restricciones $\to$ región factible $\to$ vértices $\to$ evaluación $\to$ respuesta.}

\ex{2} Un ganadero quiere mezclar dos piensos, P y Q. Cada kg de P aporta 2 unidades de vitamina A y 1 de vitamina B, y cada kg de Q, 1 unidad de A y 3 de B. Cada día el ganado necesita al menos 8 unidades de A y 9 de B. Si el kg de P cuesta 3 € y el de Q, 2 €, ¿cuántos kg de cada pienso debe comprar para que el coste sea mínimo?

\ex{2} Una persona dispone de 30\,000 € para invertir en dos fondos, A y B, con rentabilidades anuales del 4\,\% y del 6\,\%. Quiere poner al menos 10\,000 € en A y no quiere invertir en B (más arriesgado) más de lo que invierte en A. ¿Cómo debe repartir el dinero para obtener el máximo rendimiento?

\ex{3} Una empresa tiene dos almacenes: M1, con 30 toneladas de producto, y M2, con 25 toneladas. Debe servir 20 toneladas a la tienda T1 y 25 a la tienda T2. El coste de transporte por tonelada (en euros) es:

\begin{center}\small\begin{tabular}{c|cc}& T1 & T2\\\hline M1 & 5 & 3\\ M2 & 2 & 4\end{tabular}\end{center}¿Cómo se debe organizar el transporte para que el coste sea mínimo? ¿Cuál es ese coste?

\ex{2} \pau{C. Valenciana, julio 2025} Una agencia de viajes organiza excursiones a la montaña y a la playa. La agencia obtiene 700 euros de beneficio por cada excursión a la montaña y 500 euros por cada excursión a la playa. Dispone de un total de 10 autobuses y 8 guías turísticos. Cada excursión a la montaña requiere 2 autobuses y 2 guías, mientras que cada excursión a la playa requiere 2 autobuses y 1 guía.\begin{tasks}(1)
  \task ¿Cuántas excursiones a la montaña y cuántas a la playa debe organizar para obtener el máximo beneficio posible?
  \task ¿Cuál es ese beneficio máximo?
\end{tasks}

\ex{2} \pau{C. Valenciana, junio 2025} Un creador de contenido desea invertir en publicidad hasta 10\,000 euros. Ha considerado dos empresas, A y B. La empresa A garantiza un incremento anual en número de seguidores del 10\,\% de lo que se invierte en ella, y la B, del 7\,\%. El creador quiere invertir al menos 2000 euros en la empresa B y un máximo de 6000 euros en la A, invirtiendo siempre en la A al menos lo que invierte en la B.\begin{tasks}(1)
  \task ¿Cuánto dinero debe asignar a cada empresa para que el incremento de seguidores sea máximo?
  \task ¿Cuál es ese incremento máximo?
\end{tasks}

\ex{2} \pau{C. Valenciana, junio 2023} El veterinario me ha recomendado que mi perro tome diariamente un mínimo de 8 unidades de hidratos de carbono, 46 de proteínas y 12 de grasas. Una lata de la marca A contiene 4 unidades de hidratos, 6 de proteínas y 1 de grasas; una lata de la marca B, 2 unidades de hidratos, 20 de proteínas y 12 de grasas. La lata A cuesta 10 euros y la B, 16 euros.\begin{tasks}(1)
  \task ¿Cómo tendré que combinar las dos marcas para obtener la dieta deseada al precio mínimo?
  \task ¿Cuál es el precio mínimo que tendré que pagar?
\end{tasks}

\ex{2} \pau{C. Valenciana, junio 2026} Una granja quiere crear una mezcla fertilizante con dos compuestos, NutriMax y BioGrow. Se deben aplicar como mínimo 8 unidades de nitrógeno, 6 de fósforo y 10 de potasio. Un litro de NutriMax contiene 2 unidades de nitrógeno, 1 de fósforo y 1 de potasio, y un litro de BioGrow, 1 de nitrógeno, 1 de fósforo y 2 de potasio. El litro de NutriMax cuesta 5 euros y el de BioGrow, 4 euros.\begin{tasks}(1)
  \task ¿Cuántos litros de cada producto se deben utilizar para que el coste de la mezcla sea mínimo?
  \task ¿Cuál es ese coste mínimo?
\end{tasks}


% =====================================================
\seccio{5}{Casos especiales: infinitas soluciones y problemas sin solución}

\ex{2} En la carpintería del ejercicio 16 cambian los precios y ahora el beneficio es de 30 € por mesa y 10 € por silla. ¿Qué producción da el máximo beneficio? ¿Es única?\recorda{Si el óptimo se alcanza en dos vértices, se alcanza en todos los puntos del segmento que los une: ocurre cuando la función objetivo es paralela a un lado de la región.}

\ex{2} En la región del ejercicio 17 (la de los piensos), estudia si existen el máximo y el mínimo de las funciones:\begin{tasks}(2)
  \task $f(x,y)=x+y$
  \task $g(x,y)=2x+y$
\end{tasks}

\ex{2} Una empresa debe producir $x$ unidades de un producto A e $y$ de un producto B. Por capacidad, no puede fabricar más de 10 unidades en total, pero un contrato la obliga a cumplir $2x+3y\geq40$. Maximiza el beneficio $B(x,y)=3x+2y$.

\ex{3} En la región del ejercicio 6a), halla los valores de $k$ para que la función $f(x,y)=kx+y$ alcance el máximo en infinitos puntos.

\ex{2} \pau{C. Valenciana, junio 1998} Considera el triángulo de vértices $(0,\,0)$, $(2,\,8)$ y $(10,\,3)$. Determina razonadamente:\begin{tasks}(1)
  \task El punto del triángulo en el que la función $f(x,y)=-4x+y+9$ toma el máximo.
  \task El punto del triángulo en el que la función $g(x,y)=4x+y+12$ toma el máximo.
\end{tasks}

\ex{1} Razona si son verdaderas o falsas las afirmaciones siguientes:\begin{tasks}(1)
  \task Si la región factible es acotada y no vacía, la función objetivo siempre tiene máximo y mínimo.
  \task Si la región factible no es acotada, la función objetivo no tiene nunca mínimo.
  \task El óptimo de una función objetivo no constante puede estar en un punto del interior de la región.
  \task Si dos vértices dan el mismo valor máximo, el problema tiene infinitas soluciones.
\end{tasks}

\vfill
\begin{center}\color{muted}\small\faCheckCircle[regular]\ \ ¡Ánimo! Comprueba tus resultados con el resolutor de los apuntes interactivos de la unidad.\end{center}

\end{document}
